\documentclass[10pt,a4paper]{amsart}
\usepackage[margin=19mm]{geometry}
\usepackage{amsmath,amssymb,mathtools}
\usepackage[scr=rsfs,cal=euler]{mathalfa}
\usepackage{xfrac}
\usepackage[unicode,hidelinks,bookmarks=false]{hyperref}
\newcommand{\ie}{i.e.\ }
\newcommand{\cf}{cf.\ }
\newcommand{\resp}{respectively\ }
\newcommand{\Subsection}{\subsection}
\providecommand{\nobreakdash}{\nobreak\hbox{-}}
\setlength{\emergencystretch}{2em}
\multlinegap0pt
\usepackage{bm}
\usepackage{bbm}
\usepackage{dsfont}
\usepackage{xspace}
\usepackage{cancel}
\usepackage{mathtools}
\usepackage[shortlabels]{enumitem}
\usepackage{amsthm}
\makeatletter
\@ifundefined{theorem}{\newtheorem{theorem}{Theorem}}{}
\@ifundefined{assumption}{\newtheorem{assumption}[theorem]{Assumption}}{}
\@ifundefined{lemma}{\newtheorem{lemma}[theorem]{Lemma}}{}
\@ifundefined{proposition}{\newtheorem{proposition}[theorem]{Proposition}}{}
\makeatother
\newcommand{\Ctr}{\mathsf C}
\newcommand{\Ret}{\mathsf R}
\newcommand{\cM}{\mathcal M}
\newcommand{\dd}{\,\mathrm d}
\newcommand{\eps}{\varepsilon}
\newcommand{\ii}{\mathrm i}
\title[Random Tensor Estimates and Deterministic Diagonal Resolutio]{Random Tensor Estimates and Deterministic Diagonal Resolutions for Mixed Paracontrolled Operators}
\author{Guangqian Zhao}
\date{Source: arXiv:2606.16662v2 (CC BY 4.0)}
\begin{document}
\maketitle

\noindent\emph{Source and licence.} Random Tensor Estimates and Deterministic Diagonal Resolutions for Mixed Paracontrolled Operators. Version arXiv:2606.16662v2 (CC BY 4.0), distributed under the Creative Commons Attribution 4.0 International licence, as recorded in the arXiv licence metadata for this exact version and shown on the abstract page https://arxiv.org/abs/2606.16662v2. Full rights notice: licenses/arxiv-2606.16662-CC-BY-4.0.txt.\par\smallskip

\noindent\textit{Editorial scope.} Counted unit: the Proposition whose source label is \texttt{prop:diagonal} in the article (source0.tex lines 775--777, source label \texttt{prop:diagonal}), delivered together with the article's own complete original proof of it (source0.tex lines 779--876, one \texttt{proof} environment of 98 lines). No mathematics is written, repaired or re-derived: statement and proof are the article's own text line for line. Required context retained uncounted: ass:diag (lines 464--505); eq:phase-lower (lines 483--485); eq:amp-bound (lines 487--490); eq:osc-gain (lines 492--494); eq:retained-main-bound (lines 628--632); lem:dyadic-tail-assembly (lines 1430--1458). Every definition, display and label of those lines is retained unchanged. Omitted from this package and not redistributed: 1--463, 486--486, 491--491, 495--627, 633--774, 778--778, 877--1429, 1459--2044. Nothing inside a retained excerpt is omitted or altered: every retained body is delivered exactly as the article's lines are written, so no omission receipt of any kind accompanies this package. Citations: the retained excerpts carry no citation command at all, so no bibliography entry of the article is transcribed and no reference is redistributed. Reference adaptation: none was needed and none was made; every reference inside the delivered unit resolves inside it, so no unresolved reference or citation is produced. Printed slips kept literal: no printed word, symbol, hypothesis or number of the counted statement or of its proof was altered. Layout and class: the article's own class is \texttt{amsart} with packages including geometry, cmap, fontenc, inputenc, lmodern, microtype, amsmath, amssymb, amsthm, mathrsfs, enumitem, booktabs, cite, hyperref; the wrapper is the lane's pinned \texttt{amsart} editorial wrapper at 10pt on A4, which restates the theorem-like environments the retained text needs and adds the article's own macro definitions that the retained text needs (\texttt{\textbackslash Ctr}, \texttt{\textbackslash Ret}, \texttt{\textbackslash cM}, \texttt{\textbackslash dd}, \texttt{\textbackslash eps}, \texttt{\textbackslash ii}), with extra wrapper packages (bm, bbm, dsfont, xspace, cancel, mathtools, enumitem). The delivered numbering is the unit's own. Assets: no retained span contains a figure environment, a graphics-inclusion command, a PDF-page inclusion, a table or a TikZ picture; no image file is copied into this package, the package declares no asset dependency and no image file is redistributed with it. Licence: version arXiv:2606.16662v2 of this article is distributed under the Creative Commons Attribution 4.0 International licence, as recorded in the arXiv licence metadata for this exact version and shown on the abstract page https://arxiv.org/abs/2606.16662v2. A full rights notice is delivered at licenses/arxiv-2606.16662-CC-BY-4.0.txt. Characters: every retained excerpt is pure ASCII, so the delivered engine, XeLaTeX with its default Latin Modern text font, reports no missing character.
\begin{assumption}[Raw Volterra contraction criterion]\label{ass:diag}
Consider a mixed triple \(\tau=(i;j,k)\) whose covariance channel is nonzero.  There exists a number
\begin{equation}\label{eq:diag-gain}
        \Delta_{\tau}>20\eps
\end{equation}
such that every dyadic diagonal summand
\begin{equation}\label{eq:diag-summand}
        K_i(t-s,q+\ell)R_{jk}(\ell)\sigma_{jk}(\ell;s,t),
        \qquad |\ell|\sim N,\quad |q|\le c_{\rm ap}N,
\end{equation}
can be decomposed into an oscillatory part and an absolutely summable remainder as follows.

\begin{enumerate}[label=(D\arabic*),leftmargin=2.2em]
\item There is a finite index set \(\mathfrak V_\tau\), independent of \(N,Q,\Lambda\), such that the oscillatory part is a finite sum
\begin{equation}\label{eq:osc-decomp}
        \sum_{\nu\in\mathfrak V_\tau}
        e^{\ii(t-s)\Phi_\nu(q,\ell)}a_\nu(q,\ell;t,s).
\end{equation}
For some \(\vartheta_\nu\ge0\),
\begin{equation}\label{eq:phase-lower}
        |\Phi_\nu(q,\ell)|\ge c_\Phi N^{\vartheta_\nu}
\end{equation}
whenever the branch is integrated by parts, and
\begin{equation}\label{eq:amp-bound}
        |a_\nu(q,\ell;t,s)|+|\partial_sa_\nu(q,\ell;t,s)|
        \le C_\Phi N^{-\Gamma_\tau-\kappa_{jk}}
\end{equation}
for \(0\le s\le t\le T_0\).  The loop gain satisfies
\begin{equation}\label{eq:osc-gain}
        \kappa_{jk}+\vartheta_\nu+\Gamma_\tau-d\ge \Delta_\tau .
\end{equation}
If \(\vartheta_\nu=0\), no phase denominator is used; then the branch is included only through the summability in \eqref{eq:osc-gain}.

\item The remainder \(e_{N,Q}(q,\ell;s,t)\) satisfies the loop bound
\begin{equation}\label{eq:remainder-loop}
        \sup_{0\le s\le t\le T_0}\sup_{|q|\le c_{\rm ap}N}
        \sum_{|\ell|\sim N}|e_{N,Q}(q,\ell;s,t)|
        \le C_\Phi N^{-\Delta_\tau}.
\end{equation}
\end{enumerate}

\end{assumption}

\begin{equation}\label{eq:retained-main-bound}
        \|\Ret^\tau w\|_{C_TH^{s-\lambda_i}}
        +\|\Ret^\tau w\|_{L_T^1B^{\sigma-\lambda_i}_{2,\infty}}
        \le C\|w\|_{E_{T,\cM}^{2,\sigma}} .
\end{equation}

\begin{lemma}[Dyadic finite-set/tail assembly]\label{lem:dyadic-tail-assembly}
Let \(\mathfrak D\) be the countable set of dyadic triples \(\mathbf d=(N,Q,M)\) satisfying \(Q\le c_{\rm ap}N\) and \(M\lesssim N\).  Let \(X\) and \(Y\) be Banach spaces, and let \(U_{\Lambda,\mathbf d}:X\to Y\) be finite-rank operators.  Assume the following two conditions.
\begin{enumerate}[label=(\roman*)]
\item For each fixed \(\mathbf d\), the family \(U_{\Lambda,\mathbf d}\) is Cauchy in \(\mathcal L(X,Y)\) as \(\Lambda\to\infty\).
\item There exist deterministic numbers \(a_{\mathbf d}\ge0\) with \(\sum_{\mathbf d\in\mathfrak D}a_{\mathbf d}<\infty\) such that
\[
        \sup_{\Lambda,\Lambda'}
        \|U_{\Lambda,\mathbf d}-U_{\Lambda',\mathbf d}\|_{\mathcal L(X,Y)}
        \le a_{\mathbf d}
        \qquad \text{for every }\mathbf d\in\mathfrak D .
\]
\end{enumerate}
Then \(\sum_{\mathbf d\in\mathfrak D}U_{\Lambda,\mathbf d}\) is Cauchy in \(\mathcal L(X,Y)\), and the limit is independent of the order in which finite dyadic sets exhaust \(\mathfrak D\).

For a \(B_{2,\infty}^{\zeta}\)-target, the same conclusion follows from output-shell majorants.  Suppose that the contribution of all triples with output shell \(M\) is bounded by \(b_M\ge0\), with \(\sum_Mb_M<\infty\).  Then, for \(\|x\|_X\le1\),
\begin{equation}\label{eq:Besov-tail-assembly}
        \sup_{M\ge M_0} M^\zeta
        \left\|P_M\sum_{(N,Q):(N,Q,M)\in\mathfrak D} U_{\Lambda,N,Q,M}x\right\|_{L_T^1L^2}
        \le \sum_{M\ge M_0} b_M
        \longrightarrow0 .
\end{equation}
The same implications hold in \(L^p(\Omega)\) if \(a_{\mathbf d}\) and \(b_M\) dominate the corresponding \(L^p(\Omega)\)-norms of the block operator differences.  If there are random majorants \(A_{\mathbf d}(\omega)\) with \(\sum_{\mathbf d}A_{\mathbf d}(\omega)<\infty\) almost surely and
\[
        \sup_{\Lambda,\Lambda'}
        \|U_{\Lambda,\mathbf d}(\omega)-U_{\Lambda',\mathbf d}(\omega)\|_{\mathcal L(X,Y)}
        \le A_{\mathbf d}(\omega),
\]
then the same finite-set/tail argument gives almost sure operator-norm Cauchy convergence.  In particular, summable \(L^1(\Omega)\) majorants imply this pathwise hypothesis by Fubini.
\end{lemma}

\begin{proposition}[Volterra criterion for an admissible diagonal]\label{prop:diagonal}
Assume Assumption~\ref{ass:diag} for \(\tau=(i;j,k)\), and let \(s<\Delta_\tau-10\eps\).  Then the choice \(\Ret^\tau_\Lambda=D^\tau_\Lambda\), \(\Ctr^\tau_\Lambda=0\) is a dyadically admissible deterministic diagonal resolution for this label with any retained gain \(\delta_\tau<\Delta_\tau\).  The cutoff diagonals converge to a deterministic operator \(D^\tau\), and \eqref{eq:retained-main-bound} holds with \(\Ret^\tau=D^\tau\).
\end{proposition}

\begin{proof}
We suppress the mixed label and write \(\Delta=\Delta_\tau\), \(\lambda=\lambda_i\), and \(\Gamma=\Gamma_\tau\).

For the absolutely summable remainder in Assumption~\ref{ass:diag}, the diagonal kernel is already bounded after the high loop:
\begin{equation}\label{eq:D-rem-mult}
        \sup_{q,t,s}\left|\sum_{|\ell|\sim N}e_{N,Q}(q,\ell;s,t)\right|
        \lesssim N^{-\Delta}.
\end{equation}
It therefore acts on the input mode \(q\) by a scalar multiplier of size \(N^{-\Delta}\).  We next treat one oscillatory summand; the number of such summands is fixed independently of the dyadic scales.  For fixed \(q,\ell,t\), set \(a_s=a(q,\ell;t,s)\) and \(\Phi=\Phi(q,\ell)\).  The contribution to the Volterra operator is
\begin{equation}\label{eq:D-osc-one}
        \mathcal V_{q,\ell}(t)
        =\int_0^t e^{\ii(t-s)\Phi}a_s\widehat w(q,s)\dd s .
\end{equation}
If the branch has \(\vartheta>0\), integration by parts in \(s\) gives the identity
\begin{equation}\label{eq:D-IBP-full}
\begin{aligned}
        \mathcal V_{q,\ell}(t)
        =&\ \frac{e^{\ii t\Phi}a(q,\ell;t,0)\widehat w(q,0)
        -a(q,\ell;t,t)\widehat w(q,t)}{\ii\Phi} \\
        &+\int_0^t\frac{e^{\ii(t-s)\Phi}}{\ii\Phi}
        \bigl(\partial_sa(q,\ell;t,s)\widehat w(q,s)
        +a(q,\ell;t,s)\partial_s\widehat w(q,s)\bigr)\dd s .
\end{aligned}
\end{equation}
Here \(t\) is treated as an external parameter; no derivative in \(t\) is used.  By \eqref{eq:amp-bound} and \eqref{eq:phase-lower}, the summand is bounded by
\begin{equation}\label{eq:D-one-summand}
        N^{-\Gamma-\kappa_{jk}-\vartheta}
        \Bigl(\|P_Qw\|_{L_T^\infty L^2}
        +\|P_Q\partial_tw\|_{L_T^\infty L^2}\Bigr)
\end{equation}
at the level of the scalar multiplier acting on the \(q\)-shell.  Summing the high loop, which has \(O(N^d)\) lattice points, gives
\begin{equation}\label{eq:D-shell-size-refined}
        N^dN^{-\Gamma-\kappa_{jk}-\vartheta}
        \lesssim N^{-\Delta},
\end{equation}
by \eqref{eq:osc-gain}.  If \(\vartheta=0\), no integration by parts is used.  Instead,
\[
        \sum_{|\ell|\sim N}|a(q,\ell;t,s)|
        \lesssim N^{d-\Gamma-\kappa_{jk}}
        \le N^{-\Delta}
\]
directly from \eqref{eq:amp-bound} and \eqref{eq:osc-gain}.  Consequently,
for every dyadic diagonal block,
\begin{equation}\label{eq:D-dyadic-L2-refined}
        \|D_{N,Q}P_Qw\|_{C_TL^2}
        \lesssim
        N^{-\Delta}\|P_Qw\|_{L_T^\infty L^2}
        +N^{-\Delta}\|P_Q\partial_tw\|_{L_T^\infty L^2},
\end{equation}
with the derivative term omitted for purely absolute branches.

We now sum this estimate.  Because the diagonal branch is supported by \(n=q\), the output frequency is \(Q\), not \(N\).  For the Sobolev target,
\begin{equation}\label{eq:D-H-weighted}
\begin{aligned}
        \|D_{N,Q}P_Qw\|_{C_TH^{s-\lambda}}
        &\lesssim Q^{s-\lambda}N^{-\Delta}\|P_Qw\|_{L_T^\infty L^2} \\
        &\quad +Q^{s-\lambda}N^{-\Delta}\|P_Q\partial_tw\|_{L_T^\infty L^2}.
\end{aligned}
\end{equation}
By Cauchy--Schwarz in the dyadic \(Q\)-sum, with the borderline logarithm
absorbed into \(N^\eps\),
\begin{equation}\label{eq:D-first-sum}
        \sum_{Q\le c_{\rm ap}N}Q^{s-\lambda}\|P_Qw\|_{L^2}
        \lesssim_\eps N^{(s-\lambda)_++\eps}\|w\|_{L^2}.
\end{equation}
The corresponding high-shell factor is \(N^{-\Delta+(s-\lambda)_++\eps}\), which is summable under \(s<\Delta-10\eps\) because either \(s\le\lambda\), in which case \((s-\lambda)_+=0\), or \(s>\lambda\), in which case \((s-\lambda)_+<s<\Delta\).  The derivative term uses
\begin{equation}\label{eq:D-dt-weighted}
        \|P_Q\partial_tw\|_{L_T^\infty L^2}
        \le Q^{\lambda+\eps}\|\partial_tw\|_{C_TH^{-\lambda-\eps}},
\end{equation}
so that
\begin{equation}\label{eq:D-der-sum}
        \sum_{Q\le c_{\rm ap}N}Q^{s-\lambda}N^{-\Delta}
        \|P_Q\partial_tw\|_{L_T^\infty L^2}
        \lesssim N^{-\Delta+(s)_++2\eps}
        \|\partial_tw\|_{C_TH^{-\lambda-\eps}}.
\end{equation}
The high-shell sum is finite because \(s<\Delta-10\eps\); if \(s<0\), the positive part is absent.  This proves the \(C_TH^{s-\lambda}\) diagonal bound.

For the Besov target, fix the output/input block \(Q\).  From \eqref{eq:D-dyadic-L2-refined},
\begin{equation}\label{eq:D-Besov-refined}
\begin{aligned}
        Q^{\sigma-\lambda}\|D_{N,Q}P_Qw\|_{L_T^\infty L^2}
        &\lesssim N^{-\Delta}Q^{-\lambda}
        \|w\|_{L_T^\infty B^\sigma_{2,\infty}} \\
        &\quad +N^{-\Delta}
        \|\partial_tw\|_{L_T^\infty B^{\sigma-\lambda}_{2,\infty}}.
\end{aligned}
\end{equation}
The sum over \(N\gtrsim Q\) is finite because \(\Delta>0\), uniformly in \(Q\).  Hence
\[
        \|D w\|_{L_T^\infty B^{\sigma-\lambda}_{2,\infty}}
        \lesssim \|w\|_{E_{T,\cM}^{2,\sigma}}.
\]
The \(L_T^1\) estimate follows immediately with the factor \(T\).

For the cutoff limit, fix a dyadic pair \((N,Q)\).  The diagonal matrix is finite and its entries converge because the cutoff symbols converge pointwise.  The estimates above give a summable majorant independent of the cutoff.  Lemma~\ref{lem:dyadic-tail-assembly} therefore proves convergence in the same operator norms.  This proves the proposition.
\end{proof}

\end{document}
